\section{距离怎样产生开集与闭集}
\label{sec:12-Real-Analysis-Support/03-Topology-Basics/01}

回头看微积分里所有 \(\varepsilon\)--\(\delta\) 论证，它们只用到了 \(|x-y|\) 的三条运算性质：距离为零当且仅当两点重合；交换两点顺序距离不变；绕经第三点不会更近。把这三条提炼为公理，关于收敛与连续的整套推理就不再绑定实数轴——可以原封不动地搬到向量、函数、概率分布和任何你定义得出``距离''的对象上。

这套公理系统叫度量空间。

\subsection{度量空间：三条公理接管全部推理}

度量空间 \((X,d)\) 由一个集合 \(X\) 和一个距离函数

\[
d:X\times X\to[0,\infty)
\]

组成，满足三条：

\begin{enumerate}
\tightlist
\item \(d(x,y)=0\) 当且仅当 \(x=y\)（不同点之间必有正距离）；
\item \(d(x,y)=d(y,x)\)（对称）；
\item \(d(x,z)\le d(x,y)+d(y,z)\)（三角不等式——``绕路不抄近道''）。
\end{enumerate}

非负性可以从这三条推出：\(2d(x,y)=d(x,y)+d(y,x)\ge d(x,x)=0\)，不过通常还是显式列出来，便于直接使用。

三条公理没有一条提到坐标、基底或算符。它们只规定集合上任意两点之间怎么定义远近，不限制集合的元素是数还是函数。

\subsection{同一集合，不同距离：选择度量就是选择``接近''的含义}

\(\R^n\) 上至少有三种常用距离：

\[
d_p(x,y)=\norm{ x-y}_p,\qquad p=1,2,\infty.
\]

\(p=1\) 是曼哈顿距离——坐标差绝对值之和；\(p=2\) 是欧氏距离——直线段的长度；\(p=\infty\) 是最大值距离——只看偏离最大的那个坐标。在平面上，\(L_1\) 的单位球是菱形，\(L_2\) 是圆，\(L_\infty\) 是正方形。形状不同，但在有限维上它们产生同一套开集和同一套收敛概念。因为所有范数彼此等价：存在常数 \(c,C>0\) 使

\[
c\norm{ v}_p \le \norm{ v}_q \le C\norm{ v}_p \quad\text{对所有 }v\in\R^n.
\]

于是一个范数下的球总包含另一个范数下某个半径的球，也被那个范数的某个球包含。两种距离判定``足够近''的方式在极限意义上一致。

离散度量则走另一个极端：

\[
d(x,y)=
\begin{cases}
0, & x=y,\\
1, & x\ne y.
\end{cases}
\]

在这种距离下，每个单点集都是开集——取半径 \(1/2\) 的球，里面只有自己。收敛的序列必须是最终恒等的常数列，因为没有比 1 更小的正距离可用。

函数空间中最常用的距离是 sup 距离：

\[
d_\infty(f,g) = \sup_x|f(x)-g(x)|.
\]

它对应一致收敛。与之对照，\(L_1\) 或 \(L_2\) 距离则允许函数在一个小测度区域上差很多，只要整段上的积分/平方积分小就行。

选择哪种度量不取决于数学上的``对错''——所有公理都满足——而取决于你想把哪种差异判定为``显著''。

举一个非数学的例子把这点钉牢。你有两份洗澡时的水温记录：目标 40\,℃，实际水温随时间波动。比较实际曲线与目标曲线的差距，\(L_2\) 距离小意味着整段时间里的平方温差累计值小，但允许在某一秒出现一个很烫的尖峰——只要它足够短。sup 距离小则要求每一秒的最大温差都受控。前者衡量整体舒适度，后者衡量最坏瞬间的安全性。哪种更接近你心里``接近''的含义，取决于你要回答什么问题。

同样的分歧在工程中反复出现。比较两条传感器响应曲线，\(L_2\) 距离小仍允许在一个很窄的温区内出现明显偏差。若该温区恰好是过温保护的触发点，这个小小的偏差就是不可接受的——即使它在积分意义上被整个温区稀释了。

有限维中范数等价保证了同一种收敛拓扑，但这不意味着给定一个具体误差阈值时，它们表达同一项现实要求。等价常数也会随维数变化——在高维空间里，\(L_1\) 和 \(L_\infty\) 之间的常数差距会放大。

\subsection{开球与开集：每一点都有活动空间}

以 \(x\) 为中心、半径 \(r>0\) 的开球为

\[
B_r(x) = \{\,y\in X : d(x,y) < r\,\}.
\]

注意严格小于——球的``外壳''不在球内。集合 \(U\subseteq X\) 称为开集，如果对 \(U\) 中的每一点 \(x\)，都存在一个（可能依赖于 \(x\) 的）半径 \(r_x>0\)，使得

\[
B_{r_x}(x) \subseteq U.
\]

翻译成人话：开集里的每一点都带有一小块全部属于该集的私人空间。无论你站在这块领地的哪里，总可以朝任意方向稍微挪一挪而不踩到外面。

任意多个开集的并集仍是开集。有限个开集的交集是开集；无限交则不一定——反例：

\[
\bigcap_{n=1}^{\infty}\left(-\frac1n,\ \frac1n\right) = \{0\},
\]

每个 \((-1/n,1/n)\) 是 \(\R\) 中的开区间，但无限交塌缩成一个单点，单点在 \(\R\) 中不是开集（任何包含 0 的开球都包含非零点）。

\subsection{闭集：把极限点留在集合里}

集合 \(F\) 称为闭集，如果它的补集 \(X\setminus F\) 是开集。在度量空间中，这等价于一个更直观的序列判据：

\begin{quote}
若 \((x_n)\subseteq F\) 且 \(x_n\to x\in X\)，则 \(x\in F\)。
\end{quote}

闭集抓住自己的极限点。\([0,1]\) 在 \(\R\) 中是闭集：任何一个收敛且项都在 \([0,1]\) 里的数列，极限不可能跑到区间外面。\((0,1]\) 不闭，因为 \(1/n\in(0,1]\) 收敛到 0，而 \(0\notin(0,1]\)。

开与闭不是非此即彼的标签。空集 \(\varnothing\) 和整个空间 \(X\) 既开又闭。在离散空间中，每一个子集都既开又闭（所有补集也是开集——事实上所有集都是开集）。很多集合既不开也不闭：\(\Q\) 在 \(\R\) 中就处于这种中间状态——任意包含有理数的开区间里都有无理数，所以 \(\Q\) 不是开集；无理数可以是一列有理数的极限，所以 \(\Q\) 也不是闭集。

闭性还依赖于环境空间。\((0,1)\) 作为 \(\R\) 的子集不闭，但在子空间 \(X=(0,1)\) 中，它就是整个空间，既是开的也是闭的。闭集与否，取决于``全集''里包含了哪些可能的极限点。

\subsection{内部、闭包与边界：三种逐层收紧的操作}

对任意集合 \(A\subseteq X\)，可以定义三个由它生成的集：

\begin{itemize}
\tightlist
\item 内部 \(A^\circ\)：包含于 \(A\) 的最大开集。等价于所有能在 \(A\) 内部自由活动的点：存在某个半径 \(r>0\) 使 \(B_r(x)\subseteq A\)。
\item 闭包 \(\overline{A}\)：包含 \(A\) 的最小闭集。也等于 \(A\) 加上 \(A\) 的全体极限点——把从 \(A\) 中取序列能逼近到的所有点都收进来。
\item 边界 \(\partial A = \overline{A}\setminus A^\circ\)：闭包中不在内部的那些点。边界的每一点的每个邻域都同时碰到 \(A\) 与 \(A\) 的补集——你站在边界上，无论往哪个方向跨一步，都可能进到集合里或踩到集合外。
\end{itemize}

以 \(A=(0,1)\subset\R\) 检验：内部就是 \((0,1)\) 本身（开区间的每一点都有小段邻居全在区间里）；闭包是 \([0,1]\)（序列 \(1/n\) 把 0 拉进闭包，\(1-1/n\) 把 1 拉进闭包）；边界是 \(\{0,1\}\)——每一个含 0 或 1 的开区间都同时碰撞 \((0,1)\) 和它的补集。

把 \(A\) 改成 \(\Q\cap[0,1]\)，闭包仍是 \([0,1]\)——任何实数都能被有理数逼近——但内部是空集：任意包含有理数的开区间里都有无理数，所以没有一个有理数拥有全部属于 \(A\) 的开球。同一个闭包，内部为空，说明 \(A\) 是``散在''的稠密子集。

这些概念在后续章节有自己的职业出路：内部区分``严格可行''与``需要碰到边界''，在优化中出现在约束规格里；闭包出现在测度论的正则性条件中；边界出现在 Green 公式和流形上的积分。

\subsection{连续性脱掉坐标：三种等价入口}

映射 \(f:(X,d_X)\to(Y,d_Y)\) 在点 \(x\) 处连续，最直接的写法沿袭 \(\varepsilon\)--\(\delta\)：

\[
\forall\varepsilon>0,\ \exists\delta>0,\ d_X(x,y)<\delta \Longrightarrow d_Y(f(x),f(y))<\varepsilon.
\]

这个定义完全不用坐标，只用到两个空间的度量。

在度量空间中有两种等价表述。序列表述：对任意满足 \(x_n\to x\) 的序列，有 \(f(x_n)\to f(x)\)。这最接近微积分中``极限与函数可交换''的直觉。开集原像表述：\(Y\) 中每个开集在 \(f\) 下的原像

\[
f^{-1}(U) = \{\,x\in X : f(x)\in U\,\}
\]

是 \(X\) 中的开集。这个表述的威力在复合时体现得最明显——两个连续映射的复合：\(g^{-1}(f^{-1}(U))\)，如果 \(f^{-1}(U)\) 开，\(g\) 连续则它的原像也开。

三种表述描述的是同一个连续性概念，入口不同，适合的场景不同。\(\varepsilon\)--\(\delta\) 适合独立做定量估计；序列表述在度量空间中最顺手——把连续性归约为数列收敛；开集原像表述适合抽象论证和拓扑推广——它甚至不需要度量，只依赖哪些集合是开的。一旦进入一般拓扑空间，最后这个入口就变成了连续性的定义本身。
